\documentclass[11pt]{article}

% ── Packages ──────────────────────────────────────────────────────────────────
\usepackage[utf8]{inputenc}
\usepackage[T1]{fontenc}
\usepackage{lmodern}
\usepackage[margin=1in]{geometry}
\usepackage{hyperref}
\usepackage{url}
\usepackage{amsmath}
\usepackage{amssymb}
\usepackage{booktabs}
\usepackage{graphicx}
\usepackage{xcolor}
\usepackage{listings}
\usepackage{microtype}
\usepackage{parskip}
\usepackage{natbib}
\usepackage{caption}
\usepackage{subcaption}
\usepackage{multirow}
\usepackage{array}

\hypersetup{
  colorlinks=true,
  linkcolor=blue!60!black,
  citecolor=blue!60!black,
  urlcolor=blue!60!black,
  pdftitle={Affective Field Theory: A Five-Layer Emotional Memory Architecture for LLM Agents},
  pdfauthor={Gianluca Mazza},
  pdfkeywords={LLM memory, affective computing, emotion, AFT, mood-congruent retrieval, reconsolidation},
  pdfsubject={AI / Affective Computing},
}

\lstset{
  basicstyle=\ttfamily\small,
  breaklines=true,
  frame=single,
  backgroundcolor=\color{gray!8},
}

% ── Metadata ──────────────────────────────────────────────────────────────────
\title{%
  \textbf{Affective Field Theory: A Five-Layer Emotional Memory\\
  Architecture for LLM Agents}
}

\author{%
  Gianluca Mazza\\
  \texttt{info@gianlucamazza.it}\\[4pt]
  \small\textit{Independent Researcher}
}

\date{May 2026}

% ── Release metadata (managed by scripts/sync_release_metadata.py) ─────────────
% Edit release.toml, then run: make sync-metadata
\newcommand{\zenodoconceptdoi}{10.5281/zenodo.19972258}% [ssot:concept_doi]
\newcommand{\zenodoversiondoi}{10.5281/zenodo.20962443}% [ssot:version_doi]
\newcommand{\repourl}{https://github.com/gianlucamazza/emotional-memory}% [ssot:repo_url]
\newcommand{\arxivid}{}% [ssot:arxiv_id]
\newcommand{\zenodoconcepturl}{https://doi.org/\zenodoconceptdoi}
\newcommand{\zenodoversionurl}{https://doi.org/\zenodoversiondoi}

% ──────────────────────────────────────────────────────────────────────────────
\begin{document}
\maketitle

\begin{abstract}
We present \emph{Affective Field Theory} (AFT), a five-layer computational model
for emotionally-grounded memory in large language model (LLM) agents. Existing LLM
memory systems retrieve by semantic similarity and recency, neglecting the
affective modulation of human memory. AFT encodes every event with a five-component
affective fingerprint---Russell circumplex core affect, affective momentum, a
slow-adapting PAD mood field, a Scherer appraisal vector, and a resonance graph---and
retrieves with a six-signal mood-adaptive scorer gated by a Pearce-Hall
prediction-error term, validated through 127 fidelity cases over 20 psychological
phenomena. Our central result is a
characterized \emph{boundary}: affect improves retrieval exactly when it discriminates
among candidate memories. Within that regime AFT reaches \texttt{top1\_accuracy}=0.53
vs.\ 0.33 for semantic-only cosine (v2, $N=200$, SBERT; $\Delta=+0.21$ [0.15, 0.27],
$p<0.001$, $d=0.49$; replicated on e5-small-v2), and a pre-registered circularity audit
shows ${\sim}62\%$ of it is affect-discriminative by construction---the advantage
concentrates there, null on the rest. Two findings hold \emph{independent} of the
curated benchmark. (i)~Asking the LLM for valence/arousal/dominance directly is
human-validated against EmoBank (valence $r=0.79$), beating the theory-driven projection
on every axis. (ii)~Appraising the query \emph{at retrieve time} recovers most of it
with no oracle ($\Delta=+0.115$, $p<0.001$; ${\sim}59\%$ recovery, ${\sim}82\%$ on the
affect-discriminative subset)---production-reachable, yet regime-bound (null on
naturalistic dialogue, $\Delta=-0.008$). The ranking edge converts end-to-end
(LLM-judged answer accuracy 0.595 vs.\ 0.440, $p<0.001$); with automatic appraisal on
affect-free data it does not beat cosine ($\Delta=-0.087$, $p=0.002$), locating the
boundary at retrieve-time state injection. AFT encodes 25$\times$ faster than Mem0
(CPU-only) with no retrieval-time LLM call. Released as the MIT-licensed
\texttt{emotional-memory} Python library.
\end{abstract}

% ──────────────────────────────────────────────────────────────────────────────
\section{Introduction}

Memory is the substrate of agency. An LLM agent that cannot remember previous
interactions is stateless; one that remembers only facts is emotionally blind.
Human memory is neither: it is profoundly shaped by the emotional context in which
experiences were encoded and in which they are retrieved.

Bower's seminal work demonstrated mood-congruent recall~\citep{bower1981mood}:
material encoded in a positive state is more accessible when retrieved in a
matching positive state. McGaugh showed that arousal modulates consolidation
strength via the amygdala, making emotionally intense memories more durable
\citep{mcgaugh2004amygdala}. Nader and colleagues demonstrated that retrieved
memories can be updated during a lability window, a process called
reconsolidation~\citep{nader2000fear}.

These findings are rarely reflected in computational memory systems for LLM agents.
Generative Agents~\citep{park2023generative} use a composite retrieval score
including recency, importance, and semantic relevance, but no affective layer.
MemGPT~\citep{packer2023memgpt} and Mem0~\citep{chhikara2025mem0} focus on
persistence and context management without emotional modulation. MemoryBank
\citep{zhong2023memorybank} introduces Ebbinghaus-inspired decay but lacks a
valence-arousal model or mood-congruent retrieval.

We address this gap with \emph{Affective Field Theory} (AFT), a five-layer model
grounded in Russell's circumplex~\citep{russell1980circumplex}, Scherer's Component
Process Model~\citep{scherer1984appraisal}, Collins and Loftus spreading
activation~\citep{collins1975spreading}, ACT-R decay~\citep{anderson1983actr},
and Pearce-Hall associative learning~\citep{pearce1980hall}.
Each component draws on established cognitive-psychological theory; the contribution
of AFT is their integration into a unified, empirically-validated memory engine
with a pre-registered evidence programme. To our knowledge, no prior open-source
system combines all five affective layers with mood-congruent retrieval and
APE-gated reconsolidation in a single production-ready implementation.

\paragraph{Contributions.}
\begin{enumerate}
  \item A five-layer theoretical model with formal definitions for each layer (\S\ref{sec:aft}).
  \item A six-signal retrieval scorer with mood-adaptive weights (\S\ref{sec:retrieval}).
  \item An APE-gated reconsolidation mechanism (\S\ref{sec:reconsolidation}).
  \item 127 parametrized fidelity test cases validating 20 psychological phenomena (\S\ref{sec:fidelity}; intra-theory coherence validation — per-layer retrieval attribution is below detectable effect size at $N=200$, see \S\ref{sec:limitations}).
  \item A new affect-labeled dataset and comparative benchmark (\S\ref{sec:benchmark}).
  \item A pre-registered circularity-audit method that quantifies, from the data alone, how much of an affect-retrieval benchmark is favorable by construction—bounding the headline effect to its discriminative regime (Addendum~U; \S\ref{sec:limitations}).
  \item A human-gold measurement showing direct-LLM valence/arousal/dominance rating is human-validated against EmoBank and beats the theory-driven appraisal projection on every axis (Addendum~V; \S\ref{sec:limitations}).
  \item A production-reachable retrieve-time query-appraisal mechanism that recovers most of the oracle advantage \emph{without} an oracle, with its regime boundary characterized on naturalistic dialogue (Addenda~T/T2A; \S\ref{sec:limitations}).
  \item An open-source Python implementation: \href{\repourl}{\texttt{gianlucamazza/emotional-memory}} (GitHub).
\end{enumerate}

% ──────────────────────────────────────────────────────────────────────────────
\section{Affective Field Theory}
\label{sec:aft}

AFT models the affective context of a memory event as a five-dimensional field
$\mathcal{F} = (c, m, \mu, \mathbf{a}, \rho)$, where each component is defined below.
Figure~\ref{fig:aft_structures} illustrates the circumplex representation and the
resonance graph structure.

\subsection{Layer 1: Core Affect}

$c = (\nu, \alpha, \delta) \in [-1,1] \times [0,1] \times [0,1]$ is a point in
the Pleasure-Arousal-Dominance (PAD) space, where $\nu$ is valence (negative
$\leftrightarrow$ positive), $\alpha$ is arousal (calm $\leftrightarrow$ excited),
and $\delta$ is dominance (helpless $\leftrightarrow$ full control)
\citep{russell1980circumplex,mehrabian1974approach}. Core affect represents the
instantaneous affective state at encoding time. The dominance axis is populated
from the appraisal layer (coping potential, Layer~4) and is a first-class
dimension of all retrieval signals (affect proximity $s_3$ uses the 3D
Euclidean distance with normalizer $\sqrt{6}$).

\subsection{Layer 2: Affective Momentum}

$m = (v, a) \in \mathbb{R}^2 \times \mathbb{R}^2$ captures the velocity and
acceleration of affective change over a three-point history. Momentum is
non-zero when the system is in the process of emotional transition. High momentum
during encoding increases the salience of the contextual change.

\subsection{Layer 3: Mood Field}

$\mu \in [-1,1]^3$ is a PAD (Pleasure-Arousal-Dominance) vector maintained as an
exponential moving average (EMA) with coefficient $\eta$:
\[
  \mu_t = (1 - \eta)\,\mu_{t-1} + \eta\,\text{pad}(c_t)
\]
The mood field changes slowly relative to core affect, implementing the
Heideggerian notion of \emph{Stimmung} (attunement) as a slow background
predisposition~\citep[§29]{heidegger1927sein}. An optional time-decay term
$\lambda_\mu$ allows mood to regress toward a neutral baseline.

\subsection{Layer 4: Appraisal Vector}

$\mathbf{a} \in [-1,1]^5$ encodes Scherer's five Stimulus Evaluation Checks (SECs):
novelty, intrinsic pleasantness, goal relevance, coping potential, and norm
compatibility~\citep{scherer1984appraisal}. A learned or rule-based appraisal
engine maps the encoding context to an appraisal vector, which is then projected
to a core affect update via a fixed linear map. Coping potential (SEC~4) maps
directly to the dominance axis of core affect ($\delta = \text{coping\_potential}
\in [0,1]$, where 0 encodes helplessness and 1 encodes full control), so that
appraisal and core affect share a unified 3D PAD representation.

\subsection{Layer 5: Resonance Graph}

$\rho = \{(s, t, w, \tau)\}$ is a set of directed weighted edges between
memories. Link types (\emph{semantic}, \emph{emotional}, \emph{temporal},
\emph{causal}, \emph{contrastive}) follow Aristotle's laws of
association~\citep{bower1981mood}. Weights are initialized at encoding from
similarity scores and updated via Hebbian strengthening on co-retrieval
\citep{hebb1949organization}.

\begin{figure}[ht]
\centering
\begin{subfigure}[b]{0.48\textwidth}
  \centering
  \includegraphics[width=\textwidth]{figures/fig1_circumplex.pdf}
  \caption{Russell circumplex: 12 sample memories plotted by valence and
  arousal. Colour encodes consolidation strength.}
  \label{fig:circumplex}
\end{subfigure}
\hfill
\begin{subfigure}[b]{0.48\textwidth}
  \centering
  \includegraphics[width=\textwidth]{figures/fig4_resonance_network.pdf}
  \caption{Resonance link graph snapshot. Edge thickness encodes link weight;
  node labels identify memories.}
  \label{fig:resonance}
\end{subfigure}
\caption{AFT representational structures: the valence-arousal circumplex (left)
and the resonance associative graph (right).}
\label{fig:aft_structures}
\end{figure}

% ──────────────────────────────────────────────────────────────────────────────
\section{System Design}
\label{sec:design}

\subsection{Encoding Pipeline}

\begin{enumerate}
  \item \textbf{Appraisal}: optionally compute $\mathbf{a}$ from content and context.
  \item \textbf{Affect update}: derive $c_t$ from $\mathbf{a}$ or carry forward $c_{t-1}$.
  \item \textbf{State update}: update $m$ and $\mu$.
  \item \textbf{Tag}: snapshot all five layers as an \texttt{EmotionalTag}.
  \item \textbf{Embed}: compute a dense vector $e \in \mathbb{R}^d$ from the text.
  \item \textbf{Store}: persist the memory.
  \item \textbf{Resonance}: build bidirectional links against existing memories.
\end{enumerate}

\paragraph{Dual-path encoding.} Following LeDoux's dual-pathway
model~\citep{ledoux1996emotional}, a fast-path variant skips appraisal at encode
time and defers it to an asynchronous \texttt{elaborate()} call, enabling
real-time latency requirements.

\subsection{Retrieval Pipeline}
\label{sec:retrieval}

Retrieval proceeds in two passes to enable spreading activation:

\paragraph{Pass 1.} Score all candidates with an empty active set:
\[
  \text{score}(q, m) = \sum_{i=1}^{6} w_i \cdot s_i(q, m)
\]
where the six signals are: (1)~semantic cosine similarity, (2)~mood congruence,
(3)~core affect proximity, (4)~momentum alignment, (5)~ACT-R power-law decay
\citep{anderson1983actr}, and (6)~resonance boost (zero in Pass 1).

\paragraph{Pass 2.} Seed the spreading activation engine with the Pass-1 top-$k$
IDs, then re-score with resonance boosts computed by
Collins-Loftus BFS~\citep{collins1975spreading}.

\paragraph{Adaptive weights.} The weight vector $w$ is a function of the current
mood field $\mu$. High positive mood increases weight on semantic similarity;
negative mood increases weight on core affect proximity.
Figure~\ref{fig:dynamics} illustrates memory decay and mood-field evolution.

\subsection{Reconsolidation}
\label{sec:reconsolidation}

On retrieval, the Pearce-Hall associative prediction error (APE) is computed:
\[
  \text{APE}(m) = \|c_t - c_{\text{expected}}\|_1
\]
where $c_{\text{expected}}$ is an EMA of past core affect values for this memory.
When $\text{APE} > \theta$, the memory enters a lability window of duration
$\Delta_w$ seconds. Any retrieval within this window triggers reconsolidation:
the tag is blended toward the current affective state, weighted by $\text{APE}$.

\begin{figure}[ht]
\centering
\begin{subfigure}[b]{0.48\textwidth}
  \centering
  \includegraphics[width=\textwidth]{figures/fig2_decay_curves.pdf}
  \caption{ACT-R power-law decay for three arousal levels (0.1, 0.5, 0.9)
  $\times$ three retrieval-count tiers (0, 5, 15 retrievals).
  High-arousal memories decay more slowly~\citep{mcgaugh2004amygdala}.}
  \label{fig:decay}
\end{subfigure}
\hfill
\begin{subfigure}[b]{0.48\textwidth}
  \centering
  \includegraphics[width=\textwidth]{figures/fig3_mood_evolution.pdf}
  \caption{PAD mood-field trajectory over a 20-turn simulated conversation.
  The EMA smoothing of the Heidegger mood field is visible.}
  \label{fig:mood}
\end{subfigure}
\caption{System dynamics: memory decay modulated by arousal (left) and
PAD mood-field evolution over time (right).}
\label{fig:dynamics}
\end{figure}

\subsection{Performance}

\begin{sloppypar}
Table~\ref{tab:perf} reports encode and retrieve latencies measured on a single
CPU core (no GPU, \texttt{InMemoryStore}).
Regenerate with \texttt{make bench-perf} and \texttt{make reproduce-paper}.
\end{sloppypar}

\begin{table}[ht]
\centering
\caption{Encode and retrieve latency (in-process, CPU only, \texttt{InMemoryStore}).}
\label{tab:perf}
\begin{tabular}{lrrrr}
\toprule
\textbf{Operation} & \textbf{Mean (ms)} & \textbf{Median (ms)} & \textbf{p95 (ms)} & \textbf{Rounds} \\
\midrule
bench\_encode\_single & 4.39 & 4.19 & 15.94 & 1641 \\
bench\_encode\_with\_resonance & 3.13 & 3.09 & 10.32 & 917 \\
bench\_encode\_no\_resonance & 4.01 & 3.71 & 12.24 & 1328 \\
bench\_encode\_scaling[10] & 6.79 & 6.29 & 54.27 & 2187 \\
bench\_encode\_scaling[100] & 2.91 & 2.83 & 7.94 & 902 \\
bench\_encode\_scaling[1000] & 5.92 & 5.80 & 10.97 & 184 \\
bench\_retrieve\_top5[100] & 1.37 & 1.29 & 2.94 & 614 \\
bench\_retrieve\_top5[1000] & 5.60 & 5.46 & 8.41 & 176 \\
bench\_retrieve\_top5[10000] & 99.23 & 97.03 & 112.00 & 10 \\
bench\_retrieve\_varying\_topk[1] & 5.72 & 5.50 & 11.54 & 196 \\
bench\_retrieve\_varying\_topk[5] & 6.42 & 6.32 & 9.79 & 172 \\
bench\_retrieve\_varying\_topk[10] & 7.33 & 7.20 & 11.89 & 107 \\
bench\_retrieve\_varying\_topk[25] & 10.06 & 9.94 & 15.69 & 107 \\
bench\_retrieve\_with\_reconsolidation & 1.87 & 1.78 & 4.24 & 472 \\
\bottomrule
\end{tabular}

\end{table}

% ──────────────────────────────────────────────────────────────────────────────
\section{Fidelity Validation}
\label{sec:fidelity}

We validate 20 psychological phenomena through 127 parametrized test cases.
\begin{sloppypar}
Table~\ref{tab:fidelity} provides a curated overview; the per-module pass/fail
breakdown is auto-generated by \texttt{make reproduce-paper} and written to
\texttt{paper/tables/table1\_fidelity.tex}. The full test suite is available
under \texttt{benchmarks/fidelity/}.
\end{sloppypar}

\begin{table}[ht]
\centering
\caption{Fidelity test summary (127 parametrized cases, 100\% pass rate).}
\label{tab:fidelity}
\begin{tabular}{lll}
\toprule
\textbf{Phenomenon} & \textbf{Reference} & \textbf{Cases} \\
\midrule
Mood-congruent recall        & Bower 1981              & 3  \\
Arousal-enhanced memory      & McGaugh 2004            & 4  \\
ACT-R power-law decay        & Anderson 1983           & 6  \\
Spacing effect               & \citet{ebbinghaus1885memory} & 4  \\
Reconsolidation window       & Nader et al.\ 2000      & 6  \\
APE-gated lability           & Pearce \& Hall 1980     & 8  \\
Hebbian strengthening        & Hebb 1949               & 4  \\
Spreading activation         & Collins \& Loftus 1975  & 6  \\
Emotional enhancement        & \citet{kensinger2004emotional}& 4  \\
Momentum alignment           & AFT Layer 2             & 4  \\
Mood adaptive weights        & Bower 1981              & 6  \\
PAD dominance                & \citet{mehrabian1974approach} & 5  \\
Dual-path encoding           & LeDoux 1996             & 6  \\
Appraisal affect mapping     & Scherer 1984            & 8  \\
Prediction error update      & Pearce \& Hall 1980     & 6  \\
Emotion categorization       & \citet{plutchik1980general} & 14 \\
Reconsolidation window state & Nader et al.\ 2000      & 8  \\
Arousal floor                & McGaugh 2004            & 4  \\
Emotional vs cosine          & Bower 1981              & 8  \\
Design gap validation        & AFT                     & 13 \\
\bottomrule
\multicolumn{2}{l}{\textbf{Total}} & \textbf{127} \\
\end{tabular}
\end{table}

% ──────────────────────────────────────────────────────────────────────────────
\section{Comparative Benchmark}
\label{sec:benchmark}

\subsection{Dataset}

We release \textbf{affect\_reference\_v1}, a dataset of 258 affect-labeled text
examples spanning all four Russell circumplex quadrants ($Q1$--$Q4$) at three
intensity tiers (strong / moderate / mild). Each example provides valence
$\nu$, arousal $\alpha$, dominance $\delta$, and a Plutchik primary emotion
label. The dataset is available at
\texttt{benchmarks/datasets/affect\_reference\_v1.jsonl}.

\subsection{Metric: Mood-Congruent Recall@$k$}

For each of four quadrant-representative queries, we encode all 258 examples,
retrieve the top-$k$, and measure the fraction of retrieved items that are
congruent with the query's quadrant (same valence sign $\times$ arousal level).
The final score is the mean recall@$k$ across all four queries.

A random baseline would score approximately 0.25 on a balanced dataset.

\subsection{Results}

All systems use the same \texttt{all-MiniLM-L6-v2} sentence embedding model
so that differences in recall reflect scoring strategy rather than embedding
quality. For each retrieve call, AFT sets its internal affect state to the
query quadrant's centroid valence/arousal, activating mood-congruent scoring.

\begin{table}[ht]
\centering
\caption{Comparative benchmark results (recall@5, $n=258$, embedder:
\texttt{all-MiniLM-L6-v2}). Regenerate with \texttt{make bench-comparative-sbert \&\& make reproduce-paper}.}
\label{tab:comparative}
\begin{tabular}{lrrrr}
\toprule
\textbf{System} & \textbf{Recall@5} & \textbf{Encode ms} & \textbf{p50 ms} & \textbf{p95 ms} \\
\midrule
aft & 0.85 & 41.39 & 53.51 & 56.22 \\
naive\_cosine & 0.8 & 31.33 & 73.79 & 75.05 \\
recency & 0.25 & 0.01 & 0.01 & 0.03 \\
\bottomrule
\end{tabular}

\end{table}

On this coarse quadrant-level probe AFT achieves recall@5\,=\,0.85 vs.\
\texttt{naive\_cosine}\,=\,0.80 ($\Delta=+0.05$, 95\% CI [--0.15, 0.30],
$p_{\text{boot}}=0.825$).
The recency baseline confirms 0.25 as the expected floor on a balanced
four-quadrant dataset. The benchmark is intentionally narrow—it measures whether
the affect fingerprint is preserved in retrieval, not whether the multi-signal
architecture outperforms a semantic-only baseline at this granularity.
The near-tie at $N=20$ reflects a \textbf{dense-embedder ceiling effect}: with only
four quadrant centroids and recall@5 over a 20-item pool, any embedder that
correctly places memories in quadrant space saturates the metric.
Advantage of the full AFT architecture is measured on the replayable multi-session
realistic recall benchmark (v2, $N=200$, 5 challenge types $\times$ 40 scenarios),
where \texttt{top1\_accuracy}\,=\,0.53 vs.\ \texttt{naive\_cosine}\,=\,0.33
($\Delta=+0.21$ [0.15,\,0.27], $p_{\text{bootstrap}}<0.001$, $d=0.49$; SBERT
bge-small-en-v1.5). The advantage replicates on e5-small-v2
($\Delta=+0.16$ [0.09,\,0.22], $p<0.001$, $d=0.31$).
See \texttt{benchmarks/realistic/results.v2.sbert.md} and
\texttt{benchmarks/realistic/results.v2.e5.md}.
The larger, harder probe discriminates where the coarse quadrant probe cannot.

\paragraph{Architecture-attribution generalization ($H_{d2}$).}
The $H_{d1}$ result (Gate~3) replicates on the harder $N=200$ benchmark:
\textit{aft\_noAppraisal}~$=0.50$ vs.\ \textit{naive\_cosine}~$=0.38$,
$\Delta=+0.125$, $p_{\text{bootstrap}}<0.001$, $d=0.286$ (SBERT, v2).
The smaller effect on v2 (vs.\ $\Delta=+0.23$ on v1) reflects the harder
benchmark; the $\Delta>0.10$ pre-registered threshold is met.\footnote{$H_{d2}$ closure:
\href{\repourl/blob/main/benchmarks/preregistration\_addendum\_hd2\_closure.md}%
{\texttt{preregistration\_addendum\_hd2\_closure.md}};
canonical JSONs:
\href{\repourl/blob/main/benchmarks/appraisal\_confound/results.hd2.sbert.json}%
{\texttt{results.hd2.sbert.json}},
\href{\repourl/blob/main/benchmarks/appraisal\_confound/results.hd2\_it.me5.json}%
{\texttt{results.hd2\_it.me5.json}},
\href{\repourl/blob/main/benchmarks/appraisal\_confound/results.hd2\_es.sbert.json}%
{\texttt{results.hd2\_es.sbert.json}},
\href{\repourl/blob/main/benchmarks/appraisal\_confound/results.hd2\_es.me5.json}%
{\texttt{results.hd2\_es.me5.json}}.}

\paragraph{Cross-language exploratory slices.}
Spanish ($N=80$, 20~scenarios, SBERT bge-small-en-v1.5) returns
$\Delta=+0.138$, $p=0.045$, $d=0.233$ — a directional positive result.
A pre-registered power top-up (Branch~C closure) extended both the Italian and
Spanish slices to $N=120$ (30~scenarios, \texttt{multilingual-e5-small}):
Italian $\Delta=+0.058$ [$-0.042$,\;$+0.158$], $p=0.276$;
Spanish $\Delta=0.000$ [$-0.100$,\;$+0.100$], $p=1.00$.
Neither reaches significance at the pre-declared power target.\footnote{Power
top-up pre-reg:
\href{\repourl/blob/main/benchmarks/preregistration\_addendum\_hd2\_powertopup.md}%
{\texttt{preregistration\_addendum\_hd2\_powertopup.md}};
closure (Branch~C):
\href{\repourl/blob/main/benchmarks/preregistration\_addendum\_hd2\_powertopup\_closure.md}%
{\texttt{preregistration\_addendum\_hd2\_powertopup\_closure.md}};
result JSONs:
\href{\repourl/blob/main/benchmarks/appraisal\_confound/results.hd2\_it.me5.v120.json}%
{\texttt{results.hd2\_it.me5.v120.json}},
\href{\repourl/blob/main/benchmarks/appraisal\_confound/results.hd2\_es.me5.v120.json}%
{\texttt{results.hd2\_es.me5.v120.json}}.}
Cross-language evidence from the IT/ES power top-up is therefore limited to a
single exploratory Spanish-SBERT result at $N=80$; the \texttt{multilingual-e5-small}
runs at the pre-declared power do not establish the effect for Italian or Spanish.
A subsequent pre-registered French slice (Hm1, Addendum~M) at the same power ($N=120$,
me5) confirms the cross-language effect; see §Cross-language scope.

\paragraph{Cross-domain ecological replication ($H_{k1}$, DailyDialog).}
To test whether the realistic-recall advantage extends to a naturalistic corpus,
we ran a pre-registered ecological replication on DailyDialog~\citep{li2017dailydialog}
(CC~BY-NC-SA 4.0; Li et al.\ 2017, IJCNLP).
120 synthetic personas were built by concatenating 4--5 real DailyDialog dialogues
(train + validation splits, emotion-density $\geq$30\%),
with 396 affect-conditioned queries generated programmatically from turn-level
Ekman emotion labels (no LLM in the loop).
Pre-registered decision rule: paired bootstrap $n=10\,000$, seed~0, Holm $m=4$,
one-tailed; PASS requires $p_{\text{holm}}<0.05$ and $\Delta>0$.
Result: AFT \texttt{top1}~$=0.212$ vs.\ naive cosine~$=0.220$
($\Delta=-0.008$, 95\% CI $[-0.056,+0.043]$, $p_{\text{holm}}=1.000$, $d=-0.015$;
\textbf{Hk1 FAIL, Branch~B}).
The \texttt{affective\_trajectory} sub-type shows a positive trend
($\Delta=+0.103$, $d=0.186$, $N=39$) but is underpowered ($p_{\text{holm}}=0.734$).%
\footnote{Hk1 pre-registration:
\href{\repourl/blob/main/benchmarks/preregistration\_addendum\_k\_dailydialog.md}%
{\texttt{preregistration\_addendum\_k\_dailydialog.md}};
closure:
\href{\repourl/blob/main/benchmarks/preregistration\_addendum\_k\_dailydialog\_closure.md}%
{\texttt{preregistration\_addendum\_k\_dailydialog\_closure.md}};
results: \href{\repourl/blob/main/benchmarks/dailydialog/results.json}%
{\texttt{benchmarks/dailydialog/results.json}}.}
The DailyDialog FAIL is consistent with the LoCoMo negative result:
AFT's retrieval advantage is \emph{regime-specific}, materialising when the
retrieval task is explicitly designed around affective arcs and momentum shifts
(realistic\_recall\_v2) rather than naturalistic short turns.

To reproduce this table: \texttt{make bench-comparative-sbert \&\& make reproduce-paper}.

\paragraph{SOTA reference (exploratory).}
We extend the comparative probe with two production LLM-backed memory
systems---Mem0~\citep{chhikara2025mem0} and LangMem (LangGraph)---using
\texttt{gpt-4.1-mini}.  Both runs are exploratory ($H_{\text{sota}}$,
$H_{\text{v2\_sota}}$; no Holm correction).

\textit{affect\_reference\_v1 (controlled quadrant probe).}%
\footnote{Closure: \href{\repourl/blob/main/benchmarks/comparative/protocol.sota.md}%
{\texttt{protocol.sota.md}}; results: \texttt{results.sota.sbert.md}.}
Mem0 achieves recall@5\,=\,1.00 [1.00, 1.00]
($\Delta=+0.20$ vs.\ cosine, $p=0.045$); LangMem 0.90 ($\Delta=+0.10$, $p=0.254$);
AFT 0.85 ($\Delta=+0.05$, $p=0.825$).
LLM-backed memory dominates on this probe.
On encode latency, AFT (45 ms/item) is $\sim$25$\times$ faster than
Mem0 (1130 ms/item) and $\sim$3$\times$ faster than LangMem (150 ms/item),
with no LLM call at runtime.

\textit{realistic\_recall\_v2 (multi-session realistic replay).}%
\footnote{Closure: \href{\repourl/blob/main/benchmarks/comparative/protocol.sota.v2.md}%
{\texttt{protocol.sota.v2.md}}; results: \texttt{results.sota.v2.sbert.md}.}
On the same v2 dataset as the headline Hd2 result,
AFT achieves top1\_accuracy\,=\,0.535 [0.465, 0.600]
($\Delta=+0.210$ [+0.155, +0.270] vs.\ cosine, $p<0.001$, $d=0.512$).
Mem0 scores 0.330 ($\Delta=+0.005$, $p=0.945$) and
LangMem 0.365 ($\Delta=+0.040$, $p=0.233$);
neither outperforms naive cosine on this benchmark.
The advantage is consistent across challenge types:
AFT leads by $+0.275$ on affective\_arc, momentum\_alignment,
and semantic\_confound (N\,=\,40 each).

\textit{Interpretation.}
The two results are complementary, not contradictory.
\texttt{affect\_reference\_v1} tests semantic quadrant recall (four centroids,
N\,=\,258 memories); LLM-extracted summaries produce highly relevant
matches.  \texttt{realistic\_recall\_v2} tests affect-modulated multi-session
retrieval; oracle mood, momentum, and resonance scoring provide the decisive
signal.  AFT's architecture-specific advantage materialises on the task it was
designed for.  The trade-off is scope: the v2 advantage holds only under
preset oracle affect labels (see \S\ref{sec:limitations} oracle-affect caveat).
Both systems produce no LLM call at test time; the latency advantage
($\sim$25$\times$ encode) is preserved across both evaluations.

% ──────────────────────────────────────────────────────────────────────────────
\section{Related Work}
\label{sec:related}

\begin{sloppypar}
We survey 29 papers from 2022--2026; a full review is available in
\texttt{docs/research/07\_related\_work.md}. We highlight the most directly
relevant here.
\end{sloppypar}

\paragraph{LLM memory architectures.}
\citet{park2023generative} combine recency, importance, and cosine retrieval
in a memory stream for autonomous agents but include no affective modulation.
\citet{packer2023memgpt} introduce a virtual OS model for large context but
focus on persistence rather than emotional depth.
\citet{fountas2024emllm} draw on the Complementary Learning Systems framework
for episodic memory segmentation, approaching human-like memory boundaries
without affective weighting.

\paragraph{Emotion-aware retrieval.}
\citet{sumida2024lufy} select history entries by emotion similarity, validating
that arousal-modulated selection improves downstream dialogue quality — an
empirical confirmation of AFT's Layer 1 design choices.
\citet{huang2024emotionalrag} augment RAG pipelines with emotional context
vectors, finding significant quality improvement in emotional support dialogues.
\citet{wang2025anna} propose a cognitive architecture with emotion-driven memory
most similar to AFT in spirit, but without Russell circumplex formalization,
Hebbian strengthening, or APE-gated reconsolidation.
\citet{mememo2026} introduce MemEmo, a recent holistic benchmark for memory
systems with emotional context, which our affect\_reference\_v1 dataset complements.

\paragraph{Reconsolidation and updating.}
\citet{wei2026fademem} model forgetting-aware memory decay in dialogue but use
heuristic schedules rather than the ACT-R power-law modulated by encoding
arousal. No paper in our survey implements APE-gated reconsolidation windows.

\paragraph{Appraisal-based architectures.}
\citet{croissant2024chain} implement a Scherer-inspired chain-of-emotion step
before LLM response generation, mapping appraisal dimensions to emotional
states to improve emotional support dialogue — the closest prior antecedent to
AFT Layer~4. Unlike AFT, that work does not maintain a persistent affective
memory state, resonance graph, or multi-session retrieval pipeline.

\paragraph{Spreading activation in retrieval.}
\citet{gutierrez2024hipporag} use Personalized PageRank over a knowledge graph
as a spreading-activation retrieval mechanism, demonstrating its efficacy over
pure cosine retrieval on multi-hop QA.
\citet{jiang2026synapse} apply semantic spreading activation to long-context
LLM memory — both are architecturally related to AFT Layer~5, but without
valence-arousal grounding or Hebbian co-retrieval strengthening.

\paragraph{Affective computing and LLM emotion understanding.}
Several recent works examine whether LLMs can model human affect.
\citet{gandhi2024humanlike} demonstrate human-like affective cognition in
foundation models; \citet{martinez2024llmvalence} show that LLMs can estimate
valence and arousal with moderate accuracy; and \citet{liu2024emollms} release
a suite of emotion-tuned LLMs and annotation tools.
\citet{zhang2024affective} survey the emerging intersection of affective computing
and LLMs. In the memory-specific domain, \citet{xu2025amem} propose A-MEM, an
agentic memory system. \citet{damllm2025} and \citet{tan2025prospective}
address dynamic affective state management and prospective/retrospective reflection
in LLM agents, respectively. None of these systems simultaneously implement the
five-layer affective fingerprint, APE-gated reconsolidation, and Hebbian resonance
that distinguish AFT.

\paragraph{AFT differentiators.}
AFT is, to our knowledge, the only open-source system that simultaneously
provides: (1)~a Russell circumplex representation driving mood-congruent
retrieval; (2)~a five-layer affective fingerprint per memory; (3)~APE-gated
reconsolidation windows; (4)~arousal-modulated ACT-R decay with an arousal
floor; and (5)~spreading activation with Hebbian link strengthening.
Individual components draw on well-established antecedents (Russell 1980;
Collins \& Loftus 1975; Scherer 1984; Bower 1981; Hebb 1949); the contribution
of AFT is their integration into a unified, empirically-validated memory engine
with a pre-registered evidence programme.

% ──────────────────────────────────────────────────────────────────────────────
\section{Limitations}
\label{sec:limitations}

\paragraph{Embedding quality.}
Affective signals complement, but cannot compensate for, low-quality embeddings.
With a random or hash-based embedder the cosine component of the retrieval score
is uninformative, suppressing the advantage of mood-congruent signals.

\paragraph{Appraisal dependency.}
The full Layer-4 pipeline requires an external LLM or rule engine for appraisal.
Without it, the AppraisalVector defaults to a zero vector and only Layers 1--3,5
are active.

\paragraph{Affect model scope.}
AFT models affect in a three-dimensional valence-arousal-dominance space
(\textsc{pad}; \citealp{mehrabian1974approach}). The dominance dimension is
populated via coping-potential appraisal (default 0.5 when no appraisal engine
is active); no study yet isolates whether dominance adds discriminative retrieval
power over the two-dimensional baseline. Richer discrete models
(e.g., Panksepp's seven primary affects \citep{panksepp1998affective}, or the
Geneva Emotion Wheel \citep{scherer2005emotions}) may better capture
fine-grained emotional distinctions.

\paragraph{Human validation: signal yes, perceived utility not yet.}
Fidelity tests validate psychological invariants from the literature against the
model's own behavior (intra-theory validation). The \emph{appraisal signal} has now
been compared against human affect labels (Addendum~S; closure:\footnote{%
\href{\repourl/blob/main/benchmarks/preregistration\_addendum\_s\_human\_gold\_appraisal\_closure.md}%
{\texttt{preregistration\_addendum\_s\_human\_gold\_appraisal\_closure.md}}.}): on
EmoBank \citep{buechel2017emobank} (human VAD, $N=300$), \texttt{LLMAppraisalEngine}
valence is human-validated (Pearson $r=0.70$ [$0.66$, $0.75$]), while arousal
($r=0.28$) and dominance ($r=0.33$) are only weak-to-moderately positive and a
$+0.15$ positive valence bias persists; the rule-based \texttt{KeywordAppraisalEngine}
is \emph{not} human-validated ($r=0.07$). This validates the first link of the
``text $\to$ affect signal $\to$ retrieval'' chain against people, but \emph{perceived
utility} of the retrieved memories by human raters (Gate~2) remains future work---a
signal-level correlation is not a substitute for recruited raters.

The weak arousal/dominance axes are largely an artifact of the SEC$\to$projection, not a
ceiling: a pre-registered comparison (Addendum~V; closure:\footnote{%
\href{\repourl/blob/main/benchmarks/preregistration\_addendum\_v\_direct\_vad\_closure.md}%
{\texttt{preregistration\_addendum\_v\_direct\_vad\_closure.md}}.}) shows that asking the LLM
for valence/arousal/dominance \emph{directly} (a custom \texttt{AppraisalSchema}, no engine
change) beats the SEC$\to$projection on every axis vs.\ EmoBank---valence $r=0.79$ (bias
${\approx}0$), arousal $r=0.58$, dominance $r=0.43$ vs.\ $0.70/0.23/0.31$. This direct-VAD
schema ships opt-in (\texttt{DIRECT\_VAD\_SCHEMA}); the Scherer SEC schema remains the default
for theory fidelity and SEC-dependent features.
Direct-VAD arousal wins on correlation but, in Addendum~V, loses on absolute error
(MAE $0.19$ vs.\ the projection's $0.11$); a pre-registered affine calibration
(Addendum~W; closure:\footnote{%
\href{\repourl/blob/main/benchmarks/preregistration\_addendum\_w\_arousal\_calibration\_closure.md}%
{\texttt{preregistration\_addendum\_w\_arousal\_calibration\_closure.md}}.}) closes that gap:
fitting \texttt{arousal\_cal}${=}0.16\,{\cdot}\,$\texttt{arousal\_direct}${+}0.45$ on held-out
EmoBank cuts arousal MAE to $0.04$ while preserving $r$, so calibrated direct-VAD
\emph{dominates} the projection on \emph{both} $r$ ($0.57$ vs.\ $0.21$) and MAE
($0.04$ vs.\ $0.11$). The large MAE drop partly reflects EmoBank's narrow arousal variance
(shrink-to-mean); the durable claim is the dual-axis dominance via $r$-preservation, with the
calibration coefficients refittable per corpus.

\paragraph{Synthetic benchmark.}
affect\_reference\_v1 was generated synthetically. A benchmark derived from
real emotional corpora (EmpatheticDialogues, GoEmotions) would provide stronger
external validity.

\paragraph{Language and embedding scope.}
The realistic-recall advantage ($\Delta=+0.21$, SBERT bge-small-en-v1.5; v2,
$N=200$) is conditional on a dense sentence embedder; under a non-semantic
(hash) embedder the gap reduces to $\Delta\approx+0.06$ and does not reach
significance\footnote{v1.4 hash-embedder pilot; full results:
\href{\repourl/blob/main/benchmarks/realistic/results.md}{\texttt{benchmarks/realistic/results.md}}.}.
Any claim about AFT's retrieval advantage should therefore be scoped to
\emph{dense embedding regimes}.

\paragraph{Cross-language scope.}
The headline advantage is established robustly in English across both embedder
families (SBERT $d=0.49$, e5-small-v2 $d=0.31$, $N=200$).
Multilingual Italian (Hd2\_IT, $N=80$)\footnote{Addendum~H companion writeup:
\href{\repourl/blob/main/docs/research/12\_multilingual\_followup.md}{\texttt{12\_multilingual\_followup.md}};
pre-registration:
\href{\repourl/blob/main/benchmarks/preregistration\_addendum\_h.md}{\texttt{preregistration\_addendum\_h.md}}.}
shows a significant \texttt{hit@k} advantage ($\Delta\approx+0.15$--$+0.16$)
across both embedder families, but \texttt{top1\_accuracy} differences are
non-significant at $N=80$.
Spanish (Hd2\_ES, SBERT, $N=80$): $\Delta=+0.138$, $p=0.045$ — a directional
positive result.
A pre-registered power top-up to $N=120$\footnote{Power top-up closure (Branch~C):
\href{\repourl/blob/main/benchmarks/preregistration\_addendum\_hd2\_powertopup\_closure.md}%
{\texttt{preregistration\_addendum\_hd2\_powertopup\_closure.md}}.}
on \texttt{multilingual-e5-small} (Italian: $\Delta=+0.058$, $p=0.276$;
Spanish: $\Delta=0.000$, $p=1.00$) does not establish the effect at declared
power for either language.
A pre-registered French slice (Hm1, Addendum~M)\footnote{FR closure:
\href{\repourl/blob/main/benchmarks/preregistration\_addendum\_m\_fr\_closure.md}%
{\texttt{preregistration\_addendum\_m\_fr\_closure.md}}.}
on \texttt{multilingual-e5-small} at the same pre-declared power ($N=120$,
30 hand-authored native scenarios, 2-session realistic-replay design)
\textbf{confirms the cross-language effect}: $\Delta=+0.18$ [+0.11,\,+0.26]
on \texttt{top1\_accuracy} ($p{<}0.0001$, $g=0.424$, $p_{\mathrm{McNemar}}{<}0.0001$).
Cross-language evidence is therefore \emph{two-positive}: Spanish-SBERT
exploratory ($N=80$) and French-me5 confirmatory ($N=120$); $H_{d2}$ remains
the open gap on \texttt{me5} at single-session granularity (Italian and Spanish
power-top-up FAIL, Addendum~Hd2).
The appraisal scenario corpora and keyword-appraisal rules remain English-only;
the LLM-backed appraisal engine handles multilingual input.

\paragraph{External-benchmark scope.}
On the LoCoMo conversational QA benchmark (Study~S1, 1986 QA pairs, scored
$N=1540$) AFT underperforms naive RAG on both pre-registered co-primary
metrics: F1 $0.168$ vs.\ $0.271$ ($H_1$ FAIL, $\Delta=-0.101$) and
judge-accuracy $0.279$ vs.\ $0.441$ ($H_2$ FAIL, $\Delta=-0.159$);
Gate~1 not met. Affective weighting does not help on factual open-domain
QA where the optimal answer is determined by literal content recall
rather than affective coherence.
A pre-registered Pareto sweep (Addendum~J, Study~S4) tested 10 weight
configurations ($N=200$ stratified subsample, seed=42, 50 QA per category):
no configuration achieved F1 $\geq$ naive RAG on any category (Hj1 FAIL).
The closest cell was W2 ($\text{sem}=0.50$, $\text{mood}=0.30$, resonance
disabled) on multi-hop ($\Delta=-0.011$).
Per-task \texttt{base\_weights} tuning is therefore a closed negative line;
future work targeting the LoCoMo gap requires architectural changes beyond
weight tuning (e.g.\ selective affect suppression or factual/affective dual-branch
retrieval).\footnote{Pareto
pre-registration and closure:
\href{\repourl/blob/main/benchmarks/preregistration\_addendum\_j.md}%
{\texttt{preregistration\_addendum\_j.md}},
\href{\repourl/blob/main/benchmarks/preregistration\_addendum\_j\_closure.md}%
{\texttt{preregistration\_addendum\_j\_closure.md}}.
A query-type classifier (\texttt{HeuristicQueryClassifier} and
\texttt{LLMQueryClassifier}) with per-type routing tables derived from the
Addendum~J closure was implemented in v0.11.0.  Addendum~L (confirmatory,
$N=1\,540$, seed~42) tests whether closed-loop routing improves aggregate F1
over the fixed W2 baseline; execution is in progress as of 2026-05-19.}

\paragraph{Component ablations.}
The pre-registered ablation programme (Addendum~E, Study~$S_3$)\footnote{Pre-registration:
\href{\repourl/blob/main/benchmarks/preregistration\_addendum\_v3.md}{\texttt{preregistration\_addendum\_v3.md}};
Addendum~F:
\href{\repourl/blob/main/benchmarks/preregistration\_addendum\_f.md}{\texttt{preregistration\_addendum\_f.md}};
$S_3$ closure (confirmatory run):
\href{\repourl/blob/main/benchmarks/preregistration\_addendum\_s3\_closure.md}%
{\texttt{preregistration\_addendum\_s3\_closure.md}};
canonical JSONs:
\href{\repourl/blob/main/benchmarks/ablation/results.v2.sbert.json}%
{\texttt{results.v2.sbert.json}},
\href{\repourl/blob/main/benchmarks/ablation/results.v2.e5.json}%
{\texttt{results.v2.e5.json}}.}
was executed at full power on the harder \texttt{realistic\_recall\_v2}
benchmark ($N=200$ queries, both SBERT bge-small-en-v1.5 and e5-small-v2).
Individual layer ablations are \emph{not} significant: removing the mood
field ($H_a$: $\Delta=0.00$, $p_{\text{boot}}=1.000$ on SBERT; $-0.015$,
$p=0.543$ on e5) or the resonance graph ($H_b$: $\Delta=+0.030$,
$p_{\text{boot}}=0.022$, $p_{\text{holm}}=0.109$ on SBERT; $+0.075$,
$p_{\text{boot}}<0.001$, $p_{\text{holm}}=0.001$ on e5 in the opposite
direction) does not reduce top1\_accuracy under family-wise correction.
Removing appraisal is invariant on this benchmark
($H_c$ PASS: $\Delta=0.00$ on SBERT, $-0.005$ on e5). The system-level
architecture advantage observed in $H_{d1}/H_{d2}$ (vs.\ \texttt{naive\_cosine},
$\Delta=+0.13$ to $+0.23$) is therefore replicable but \emph{cannot be
attributed to any single layer in isolation}: layers are jointly sufficient
and likely interact; per-layer attribution at $N=200$ is below the
detectable effect size and would require either a higher-$N$ design or
interaction-term modelling.
A symmetric dual-path encoding variant remains destructive
($H_{e1}$: $\Delta=-0.185$ on SBERT, $-0.275$ on e5;
\textit{aft\_keyword\_synchronous}: $\Delta=-0.440$ on both embedders,
replicating $H_{f1}$). Reconsolidation continues to register null
($H_{e2}$: $\Delta=+0.010$ on both embedders); this is expected because
the v2 benchmark issues each query once per persona (single-session
design), so the APE-gated lability window is never triggered.
Validating reconsolidation requires a longitudinal benchmark with
repeated retrievals of the same memory under a changing affective state.
Dual-path encoding is therefore retained as opt-in only; semantically-grounded
LLM appraisal (non-keyword) is recommended when an appraisal engine is configured.

\paragraph{Resonance magnitude amplification on e5-small-v2.}
The ablation finding for the resonance layer ($H_b$: $\Delta=+0.075$,
$p_{\text{boot}}<0.001$ on e5-small-v2) is a confirmed \emph{magnitude amplification}, not
a sign reversal: per-challenge decomposition (Addendum~I, post-hoc) shows
$\Delta \geq 0$ for every (embedder $\times$ challenge-type) cell, with the
amplification concentrated on \texttt{semantic\_confound} (e5 $\Delta=+0.125$
vs.\ SBERT $\Delta=+0.025$).
Hi3 (Addendum~I closure, $N=500$, seed$=$1, Holm $m=3$) confirms this
cross-embedder amplification at the pre-registered threshold:\footnote{%
Hi3 closure:
\href{\repourl/blob/main/benchmarks/preregistration\_addendum\_i\_closure.md}%
{\texttt{preregistration\_addendum\_i\_closure.md}}.}
\texttt{semantic\_confound} $\Delta=+0.090$ [0.030, 0.160], $d=0.257$,
Holm-adj $p=0.023$ (\textbf{PASS}); \texttt{recency\_confound} $\Delta=+0.070$,
$d=0.239$, Holm-adj $p=0.023$ (PASS); \texttt{affective\_arc} $\Delta=+0.010$,
$p=0.380$ (FAIL).
The amplification is multi-channel on semantic and recency challenges but does
not extend to affective-arc queries.
The most plausible explanation remains geometry-specific: e5-small-v2 clusters
semantically related items more tightly than SBERT-bge, causing spreading
activation to over-fire on topic-near distractors.
Link-count instrumentation (both embedders saturate the top-5 link cap,
mean~$\approx 5.0$) does not yet differentiate the mechanism at link-density
granularity; link-type and link-strength distributions remain unmeasured.
Future work: instrument per-link-type and per-link-strength statistics
across embedder families to test whether e5's tighter semantic geometry
concentrates high-strength links on topic-near distractors (Hi2 mechanism
study).

\paragraph{Oracle-affect operationalization.}
The appraisal-confound studies ($H_{d1}$, $H_{d2}$, cross-language Hd2\_IT/ES, Hm1~FR)
compare \texttt{aft\_noAppraisal}---AFT using \emph{preset} valence/arousal
values from the benchmark---against \texttt{naive\_cosine}, which has no access
to affect at all. This design isolates the \emph{retrieval architecture} given
perfect affect, but does not test the \emph{appraisal pipeline} under real
inference conditions. The Ha2/Hb2 results show that keyword appraisal
is destructive ($\Delta=-0.39$ vs.\ naive\_cosine), so the architecture
advantage under automatic appraisal was tested in Addendum~G
(executed 2026-05-07; closure:\footnote{%
\href{\repourl/blob/main/benchmarks/preregistration\_addendum\_g\_closure.md}%
{\texttt{preregistration\_addendum\_g\_closure.md}}.})
with \texttt{LLMAppraisalEngine} in dual-path mode (\texttt{gpt-5-mini})
on $N=200$ queries from \texttt{realistic\_recall\_v3\_noAF},
a dataset not pre-seeded with oracle affect values.
\textbf{Hg1 FAIL}: $\Delta=-0.010$ [$-0.055$, $0.035$], $d=-0.032$,
$p_{\text{one}}=0.367$---the architecture advantage does not extend to
LLM dual-path appraisal at this sample size.
A follow-up recalibrated the SEC$\to$valence/arousal mapping (Addendum~O;
held-out calibration PASS, cutting valence bias $+0.200\to+0.072$) and re-ran the
affect-free test on \texttt{realistic\_recall\_v4\_noAF}, a leakage-free set
disjoint from the calibration data ($N=160$; closure:\footnote{%
\href{\repourl/blob/main/benchmarks/preregistration\_addendum\_p\_hg1\_rerun\_closure.md}%
{\texttt{preregistration\_addendum\_p\_hg1\_rerun\_closure.md}}.}).
The boundary holds, more sharply: \texttt{naive\_cosine}~$=0.887$ significantly
outperforms \texttt{aft\_llm\_dual}~$=0.800$ (\textbf{Hp1 FAIL}: $\Delta=-0.087$
[$-0.144$, $-0.031$], $p=0.0018$, $d=-0.24$). Two exploratory contrasts isolate
why: dual\,$>$\,neutral ($\Delta=+0.056$, $p=0.030$) shows the LLM-inferred affect
carries genuine signal, and dual\,$>$\,sync ($\Delta=+0.512$, $d=0.95$) confirms
the deferred dual-path schedule is essential---yet the affect channel is a net
distractor when queries are affect-free and semantics already discriminate.
The oracle-affect circularity remains the operative scope boundary of the
$H_{d1}$/$H_{d2}$ claims.
A final pre-registered study (Addendum~Q; closure:\footnote{%
\href{\repourl/blob/main/benchmarks/preregistration\_addendum\_q\_affect\_gating\_closure.md}%
{\texttt{preregistration\_addendum\_q\_affect\_gating\_closure.md}}.})
tested the remaining routing synthesis---apply the affect channel only when the
query is affect-discriminative---on \texttt{realistic\_recall\_v5\_gate}
($N=200$: 100 affect-free + 100 affective queries with ground-truth gate labels,
frozen pre-run), with gating implemented as a front-router (affect-free queries
answered from a cosine index identical to the baseline; affective queries by the
full engine). \textbf{Hq2 PASS} ($\Delta=+0.080$, $p_{\text{Holm}}=0.0009$):
gating recovers the \emph{entire} always-on penalty---the gated system scores
exactly cosine on affect-free queries. But \textbf{Hq1/Hq3 FAIL}: LLM-inferred
affect loses to cosine on the affective subset itself (anchored tie-break
challenge: $0.160$ vs.\ $0.280$), and even a ground-truth-gated arm stays
significantly below cosine ($\Delta=-0.045$, $p=0.0024$)---the bottleneck is the
channel, not the gate classifier. This sharpens the oracle-affect boundary into a
\emph{state-injection} boundary: AFT's affective signals compare memories against
the runtime state and the query itself is never appraised, so the $H_{d}$-family
protocols---which inject each query's \texttt{state} before retrieval---perform an
alignment the session trajectory does not supply for free. The affect-routing
line (reshaping the channel while keeping the oracle state) is closed.

A pre-registered study then \emph{moved} this boundary (Addendum~T; closure:\footnote{%
\href{\repourl/blob/main/benchmarks/preregistration\_addendum\_t\_query\_appraisal\_closure.md}%
{\texttt{preregistration\_addendum\_t\_query\_appraisal\_closure.md}}.}). Replacing the
oracle \texttt{state} with the query's affect \emph{appraised at retrieve-time}
(direct-VAD on the query text, then injected as the state) makes AFT beat cosine
\textbf{without any oracle}: $\Delta=+0.115$ [$+0.055$, $+0.180$], $p<0.001$ on
\texttt{realistic\_recall\_v2} ($N=200$, SBERT), recovering ${\sim}59\%$ of the oracle
advantage overall and ${\sim}82\%$ on the affect-discriminative subset; the appraised
query affect tracks the oracle state (valence $r=0.80$, arousal $r=0.56$). The
affect-discriminative advantage is therefore \textbf{largely production-reachable}---the
oracle is not required---but this recovery is \emph{regime-bound}. A pre-registered
naturalistic re-test (Addendum~T2A; closure:\footnote{%
\href{\repourl/blob/main/benchmarks/preregistration\_addendum\_t2a\_naturalistic\_query\_appraisal\_closure.md}%
{\texttt{preregistration\_addendum\_t2a\_naturalistic\_query\_appraisal\_closure.md}}.}) applied
the same retrieve-time query appraisal to DailyDialog ($N=120$, 396 queries, multilingual-e5-small)
and did \emph{not} beat cosine ($\Delta=-0.008$, $p_{\mathrm{holm}}=1.000$, 0/3 directional
types), reproducing the $H_{k1}$ null---even though the query appraisal is faithful here too
(valence $r=0.69$, arousal $r=0.74$). The production-reachable recovery is thus confined to the
affect-discriminative regime and does not extend to naturalistic, affect-sparse dialogue.

\paragraph{Controlled benchmark scope.}
\texttt{realistic\_recall\_v2} was designed with 4 of 5 challenge types where
AFT's affective signals are theoretically discriminative
(\texttt{affective\_arc}, \texttt{momentum\_alignment},
\texttt{semantic\_confound}, \texttt{same\_topic\_distractor}).
Only \texttt{recency\_confound} is valence-neutral.
The aggregate advantage ($\Delta=+0.205$ SBERT; $+0.155$ e5) should therefore
be read as \emph{AFT advantage when affective context discriminates}, not as
uniform superiority across all retrieval tasks. This is consistent with two
independent negative results: the LoCoMo factual QA benchmark (Gate~1 FAIL)
and the DailyDialog ecological replication ($H_{k1}$ FAIL; see above).
In both cases, either affect is non-discriminative (LoCoMo) or the affective
content per turn is too sparse for the 6-signal scorer to outperform semantic
similarity (DailyDialog short turns, $\approx$17\% emotion-bearing turns before
density filtering). Scenarios with valence-neutral discrimination criteria are
under-represented in v2. A pre-registered circularity audit (Addendum~U; closure:\footnote{%
\href{\repourl/blob/main/benchmarks/preregistration\_addendum\_u\_circularity\_audit\_closure.md}%
{\texttt{preregistration\_addendum\_u\_circularity\_audit\_closure.md}}.}) quantifies this:
computing per query, from the data alone, whether the gold is cosine-unsolvable \emph{and}
the affect-closest candidate to the query state, \textbf{62.5\% of v2 queries are
affect-favorable by construction}. The aggregate advantage is concentrated there
($\Delta=+0.304$ [$+0.224$, $+0.384$]) and is \textbf{null on the neutral 37.5\%}
($\Delta=+0.013$ [$0.000$, $+0.040$], $p=0.63$). The $+0.205$/$+0.18$ headline figures
should thus be read as the advantage on a benchmark that is ${\sim}62\%$
affect-discriminative \emph{by construction}, not a regime-independent effect size.

\paragraph{Downstream value (oracle-affect regime).}
A pre-registered end-to-end study (Addendum~R; closure:\footnote{%
\href{\repourl/blob/main/benchmarks/preregistration\_addendum\_r\_downstream\_closure.md}%
{\texttt{preregistration\_addendum\_r\_downstream\_closure.md}}.}) tested whether the
retrieval-ranking advantage converts to \emph{answer quality} once a generator consumes
the retrieved memories. On \texttt{realistic\_recall\_v2} ($N=200$, oracle affect), an
\mbox{encode$\to$retrieve$\to$generate$\to$judge} pipeline---identical generator and
LLM judge for both systems, only the retrieval stage differing---yields
\textbf{Hr1/Hr2 PASS}: AFT LLM-judged answer accuracy $0.595$ vs.\ cosine $0.440$
($\Delta=+0.155$ [$+0.095$, $+0.220$], $p<0.001$, Holm- and McNemar-corrected),
token-F1 $\Delta=+0.152$. The $+0.205$ ranking edge converts roughly one-to-one to a
$+0.155$ answer-quality edge. This is the affect-discriminative, oracle-affect regime
and is therefore \emph{bounded by the same state-injection boundary} as the
$H_{d}$-family results; it does \emph{not} contradict the LoCoMo negative, which is the
oracle-free, naturalistic regime. The two coexist: AFT converts ranking into answers
when affect is discriminative and supplied, and loses when it is not.

% ──────────────────────────────────────────────────────────────────────────────
\section{Reproducibility}
\label{sec:reproducibility}

All experiments are reproducible via:
\begin{lstlisting}
git clone https://github.com/gianlucamazza/emotional-memory
pip install -e ".[dev,bench]"
make reproduce-paper        # generates paper/tables/
make bench-fidelity         # 127 fidelity tests
make bench-comparative      # recall@k benchmark
\end{lstlisting}

\begin{sloppypar}
The repository includes a \texttt{CITATION.cff} file.
A persistent DOI is available at \href{\zenodoconcepturl}{\texttt{\zenodoconceptdoi}}.
\end{sloppypar}
\texttt{make reproduce-paper} regenerates Tables~\ref{tab:perf} and~\ref{tab:comparative}
from live benchmark runs; Table~\ref{tab:fidelity} is a curated overview
(auto-generated module-level results at \texttt{paper/tables/table1\_fidelity.tex}).
Expected runtime on a laptop CPU: ${\sim}$5\,min for fidelity, ${\sim}$2\,min for
the comparative benchmark. Requires Python~$\geq$\,3.11.

% ──────────────────────────────────────────────────────────────────────────────
\section{Conclusion}

We presented Affective Field Theory, a five-layer model for emotionally-grounded
LLM memory, and the \texttt{emotional-memory} Python library that implements it.
The model is validated through 127 fidelity test cases covering 20 psychological
phenomena, and benchmarked against a semantic-only baseline on a new affect-labeled
dataset. The library is MIT-licensed, fully typed, and available for immediate use
in LangChain pipelines and custom LLM agents.

A staged evidence programme has produced a series of pre-registered studies
(Addenda A--V and T2A, plus the Hd2 power top-up) since initial submission. The expanded realistic-recall benchmark ($N=200$, five
challenge types, two embedder classes) confirmed a decisive advantage (SBERT:
$\Delta=+0.21$, $p<0.001$, $d=0.49$; e5-small-v2: $\Delta=+0.16$, $p<0.001$;
Gate~3 CLOSED via the appraisal-confound ablation\footnote{$H_{d1}$:
\textit{aft\_noAppraisal}~$=0.78$ vs.\ \textit{naive\_cosine}~$=0.55$,
$\Delta=+0.23$ [+0.12, +0.34], $d=0.52$, $N=100$, v1.4, seed=1 (Addendum~D;
\href{\repourl/blob/main/docs/research/audit\_2026-04.md}{\texttt{audit\_2026-04.md}}, \S{}G3).},
with $H_{d2}$ generalization confirmed: $\Delta=+0.125$ on SBERT v2 ($N=200$).
Spanish cross-language ($N=80$, SBERT, $\Delta=+0.138$, $p=0.045$) provides a
directional positive data point; a pre-registered power top-up to $N=120$ on
\texttt{multilingual-e5-small} (Italian and Spanish) did not replicate the effect
(both $p>0.25$, Branch~C closure).
A pre-registered French slice (Hm1, Addendum~M, $N=120$, me5) confirms the
cross-language effect: $\Delta=+0.18$, $p<0.0001$, $g=0.424$ (Branch~A PASS).
A pre-registered layer-ablation at power (Study~$S_3$, $N=200$, both embedder
classes) finds that per-layer contributions are not individually isolatable:
$H_a$ (no\_mood) and $H_b$ (no\_resonance) are non-significant on both
embedders, while the system-level advantage over naive cosine remains robust.
The external LoCoMo conversational
QA benchmark (1986 QA pairs, Study S1) produced a negative result: AFT
underperforms naive RAG (F1 0.168 vs.\ 0.271; Gate 1 not met), revealing that
affective weighting does not help on factual open-domain QA.
The pre-registered Pareto sweep (Study S4, Addendum~J) confirmed this scope
limitation holds across all 10 weight configurations (Hj1 FAIL): per-task
\texttt{base\_weights} tuning cannot close the gap.
A pre-registered ecological replication on DailyDialog ($H_{k1}$, $N=120$
personas, 396 affect-conditioned queries) did not replicate the advantage
(Branch~B FAIL, $\Delta=-0.008$, $d=-0.015$), confirming that the regime boundary
is naturalistic affect density rather than dataset size.
Two findings of this revision hold \emph{independent} of the curated benchmark and
are the most portable outcomes of the programme: direct-LLM VAD rating is
human-validated against EmoBank (Addendum~V, valence $r=0.79$, beating the
theory-driven projection on every axis), and retrieve-time query appraisal recovers
most of the oracle advantage with no oracle (Addendum~T, $\Delta=+0.115$, $p<0.001$,
${\sim}59\%$ recovery)---production-reachable, though a naturalistic re-test bounds it
to the affect-discriminative regime (Addendum~T2A, $\Delta=-0.008$). Per-category
weight routing, query-type classification, and per-query affect gating were tested and
are closed negatives (Addenda~J/L/Q).
Open work: a human-evaluation pilot to ground-truth affective coherence ratings
(protocol in \texttt{benchmarks/human\_eval/}); broader multilingual coverage
beyond the English/Italian/Spanish/French slices validated here;
and identifying naturalistic corpora with sufficient affect density (e.g.\
affective-trajectory proportion $>$30\%) where the DailyDialog-style replication
could be expected to pass.
The dominance axis (PAD third dimension) is now a first-class \texttt{CoreAffect}
field as of v0.8.2 (commit \texttt{8b9ddbe}); no pre-registered comparison of
3D PAD vs.\ 2D valence-arousal retrieval scoring has been executed.
Future work: a pre-registered attribution study isolating whether dominance
improves top-$k$ accuracy on emotion-rich corpora beyond the 2D baseline.
Vector-store adapters for Qdrant and ChromaDB shipped in v0.9 as optional extras
(\texttt{[qdrant]}, \texttt{[chroma]}), each implementing the duck-typed \texttt{MemoryStore}
protocol in strict vector-first mode.
Expanded baselines were executed in this revision: Mem0 and LangMem on
\texttt{affect\_reference\_v1} (LLM-backed memory dominates on the simple probe,
Mem0 recall@5\,=\,1.00 vs.\ AFT 0.85) and on \texttt{realistic\_recall\_v2}
(AFT top1\,=\,0.535 $\gg$ Mem0\,=\,0.330, LangMem\,=\,0.365; see
\S\ref{sec:benchmark} SOTA reference).
SOTA replication on the LoCoMo factual QA benchmark remains open work.

% ──────────────────────────────────────────────────────────────────────────────
\section*{Acknowledgements}
% TODO(author): add funding sources, institutional support, and individual
% contributors before camera-ready submission.
This work builds on the open-source scientific Python ecosystem; we thank the
maintainers of the libraries on which \texttt{emotional-memory} depends.

% ──────────────────────────────────────────────────────────────────────────────
\bibliographystyle{plainnat}
\bibliography{refs}

\end{document}
